\section{The standard barrier after affine restriction}\label{sec:restricted}
The rank of a certificate concerns a single support direction. A barrier
must control every direction in its domain, including directions inside a
projection fiber. We now make this distinction quantitative.

\subsection{The parameter and two boundary tests}
Write $L=X^\circ+T$ for the effective affine slice of a product of
Euclidean Jordan algebras, after minimal-face reduction, and set
$\Omega=L\cap\intr K$. ``Effective'' means that redundant variables
that do not change the cone factors have been removed. For the standard
barrier
\[
 F(X)=-\sum_i\log\det_i X_i,
\]
its least gradient parameter on this slice is
\begin{equation}\label{eq:exact-gradient-parameter}
 \nustd=\sup_{X\in\Omega}\ \sup_{0\ne H\in T}
 \frac{DF(X)[H]^2}{D^2F(X)[H,H]}.
\end{equation}
The restricted Hessian is positive definite. The standard third-derivative
self-concordance inequality is preserved under affine restriction, so
\eqref{eq:exact-gradient-parameter} is exactly the least
self-concordant-barrier parameter of this function. It is not
the infimum over all barriers on $\Omega$. The ambient Jordan rank is an
upper bound for $\nustd$ but need not equal it.

\begin{lemma}[Boundary order]\label{lem:boundary-order}
For every $\bar X\in L\cap K$,
\begin{equation}\label{eq:nullity-barrier}
 \nustd\geq\sum_i\bigl(r_i-\jrank\bar X_i\bigr).
\end{equation}
In particular, the bound holds at every point of every fiber over the
boundary, not only at a chosen boundary selection.
\end{lemma}
\begin{proof}
Let $X(t)=(1-t)\bar X+tX^\circ$, $0<t<1$, and put
$z_i=r_i-\jrank\bar X_i$. Apply the positive invertible quadratic
representation $P((X_i^\circ)^{-1/2})$. It sends $X_i^\circ$ to the unit,
preserves rank, and changes the determinant by a positive constant.
Diagonalizing the transformed $\bar X_i$ therefore gives
\[
 \det_i X_i(t)=c_i t^{z_i}\prod_{j=1}^{r_i-z_i}
 ((1-t)\lambda_{ij}+t),\qquad c_i>0,\quad\lambda_{ij}>0.
\]
If $z=\sum_i z_i$, the restriction $f(t)=F(X(t))$ satisfies
$f'=-z/t+O(1)$ and $f''=z/t^2+O(1)$. For $z>0$, the quotient
$(f')^2/f''$ tends to $z$; for $z=0$ the claim is immediate.
These statements use Jordan eigenvalues and apply also to the Albert
algebra.
\end{proof}

\begin{lemma}[Recession and complementary boundary order]\label{lem:recession}
Suppose $D\ne0$ is a recession direction of $L\cap K$. In each factor
let $c_i=\supp D_i$, $f_i=e_i-c_i$, and let
$C_i=P(f_i)\bar X_i$, where $\bar X\in L\cap K$. Define
\[
 r=\sum_i\jrank D_i,\qquad
 z=\sum_i\bigl(\rank f_i-\rank_{V_i(f_i,1)} C_i\bigr).
\]
Then $\nustd\geq r+z$.
\end{lemma}
\begin{proof}
The two-dimensional affine family
$X(\sigma,t)=(1-\sigma)\bar X+\sigma X^\circ+tD$ is interior for
$0<\sigma<1$ and $t>0$. Put $A(\sigma)=(1-\sigma)\bar X+\sigma X^\circ$.
The Jordan Schur determinant formula \citep{GS2010}, applied to the
Peirce decomposition for $c_i$, gives
\begin{equation}\label{eq:recession-leading}
 \det_i(A_i(\sigma)+tD_i)
 =t^{r_i'}\det_{c_i}D_i\,
   \det_{f_i}(P(f_i)A_i(\sigma))\,[1+O_\sigma(t^{-1})],
 \qquad r_i'=\jrank D_i.
\end{equation}
For a zero $D_i$ the formula means the unchanged determinant and
$r_i'=0$. The Schur correction is quadratic in the mixed Peirce part
and linear in the inverse of the $c_i$ part, hence is $O_\sigma(t^{-1})$.
Alternatively, the determinant is a polynomial in $(\sigma,t)$ whose
leading coefficient is the positive coefficient displayed in
\eqref{eq:recession-leading}. On every compact subinterval of
$0<\sigma<1$, the expansion may consequently be differentiated twice,
including after multiplication of each $t$ derivative by $t$.

Write $f(\sigma,t)=F(X(\sigma,t))$ and
$g(\sigma)=-\sum_i\log\det_{f_i}(P(f_i)A_i(\sigma))$, omitting rank-zero
factors. As $t\to\infty$ with $\sigma$ fixed,
\[
 tf_t\to-r,\quad t^2f_{tt}\to r,\quad
 f_\sigma\to g',\quad f_{\sigma\sigma}\to g'',\quad
 tf_{t\sigma}\to0.
\]
The compression $P(f_i)X_i^\circ$ of the Slater point is positive
definite in $V_i(f_i,1)$ whenever this algebra has positive rank. The
determinant computation of Lemma~\ref{lem:boundary-order}, applied to
$g$, gives
$\sigma g'\to-z$ and $\sigma^2g''\to z$ as $\sigma\downarrow0$.
At each point test the affine direction $H=tD+\sigma(X^\circ-\bar X)$.
Its gradient tends, in the iterated limit ($t\to\infty$ first, then
$\sigma\downarrow0$), to $-(r+z)$ and its
Hessian quadratic form to $r+z$. Equation~\eqref{eq:exact-gradient-parameter}
therefore gives $\nustd\geq r+z$. The derivative expansions are needed
here: a bounded remainder in the barrier value alone would not justify
the conclusion.
\end{proof}

\begin{proposition}[General lower bound and minimum under global selection]\label{prop:barrier-frontiers}
For a finite repeatable symmetric-cone dictionary with $B>0$, every
relative-Slater affine lift of $\B_2^s$ satisfies
\begin{equation}\label{eq:general-barrier-bracket}
 \left\lceil\frac{s-1}{B}\right\rceil\leq\nustd.
\end{equation}
There is a bounded lift with parameter $\kappa_{\mathscr K}(s)$ from
\eqref{eq:smooth-kappa}. Among lifts of $\prod_a\B_2^{s_a}$ that admit
globally bi-$C^1$ full-row selections \eqref{eq:full-rows} with genuine
certificates, the minimum of $\nustd$ is exactly
$\sum_a\kappa_{\mathscr K}(s_a)$.
For a single ball, every unbounded lift satisfies the stronger bound
$\nustd\geq1+\lceil(s-1)/B\rceil$.
\end{proposition}
\begin{proof}
The dense-open rank bound for full certificate fibers in
Theorem~\ref{thm:rank-frontier} and complementarity force a
boundary tuple with nullity at least $\lceil(s-1)/B\rceil$; apply
Lemma~\ref{lem:boundary-order}. For the upper bound use
\eqref{eq:fixed-tail}. Its determinant is $t_G-\norm{w_G}^2$, so its
barrier is
\begin{equation}\label{eq:grouped-barrier}
 F=-\sum_{G=1}^{g}\log(t_G-\norm{w_G}^2),\qquad\sum_Gt_G=1.
\end{equation}
Each summand has parameter one on its paraboloid domain. Indeed, if
$h=t-\norm w^2$, then
$D^2(-\log h)[H,H]=(Dh[H]/h)^2+2\norm{H_w}^2/h$.
The third-derivative inequality follows by restriction of the rank-two
Lorentz log determinant with one perspective coordinate fixed.
By Cauchy--Schwarz, the sum has parameter at most $g$.
At any sphere point set $t_G=\norm{w_G}^2$. Each fixed-tail block then
has nullity one, so Lemma~\ref{lem:boundary-order} gives $\nustd\geq g$,
and hence equality.
For the direct Lorentz section $t=1$ of $Q_{s+1}$, $F=-\log(1-\norm x^2)$ has parameter one:
its squared dual gradient norm is
$2\norm x^2/(1+\norm x^2)$, with supremum one.
Products give the stated upper bound, and
Theorem~\ref{thm:smooth-product-rank} gives the matching nullity lower
bound.

For the last assertion, a nonempty closed unbounded convex slice has a
nonzero recession direction $D\in K\cap T$; compactness of its projection
implies $\pi D=0$. Choose a generic support certificate $Y$ of rank at
least $q=\lceil(s-1)/B\rceil$ and any tuple $\bar X$ over that support.
The whole-slice identity gives $\ip D Y=0$, so $Y_i$ lies in the face
complementary to $\supp D_i$. Its complementarity with $\bar X_i$ implies
that the compressed nullity $z$ in Lemma~\ref{lem:recession} is at least
$q$. Since $\sum_i\jrank D_i\geq1$, that lemma proves the assertion.
\end{proof}

Consequently, when $B$ does not divide $s-1$, the minimum over all lifts
of a single ball is exact without selection hypotheses: both bounds
equal $\lceil s/B\rceil$. The same holds for Hermitian order caps with
$B=a(R-1)$, where $a=\delta=\dim_\R\mathbb F$, and for Lorentz
dimension caps with $B=d-2$.
Proposition~\ref{prop:barrier-frontiers} leaves open bounded lifts with
$s=qB+1$. Theorems~\ref{thm:one-channel} and~\ref{thm:wide-cap} settle
$q=1$ and, for Hermitian and Lorentz caps, $2\leq q\leq B+1$; for
these caps the remaining range is~\eqref{eq:unsettled-range}, and for
other dictionaries bounded lifts with $q\geq2$ remain open.
If every symmetric cone of real dimension at most $d$ is permitted, the
maximum primitive capacity is $B=d-2$, attained by $Q_d$. Hence, over
lifts with smooth global selections, the minimum standard-barrier
parameter is $\lceil s/(d-2)\rceil$ for $3\leq d<s+1$, and one for
$d\geq s+1$. For several source balls these minima add. This
statement allows all symmetric cone types; higher-rank types cannot
improve the optimum.

\subsection{An exact result for \texorpdfstring{$s=B+1$}{s=B+1} without selection hypotheses}
\begin{theorem}[One-channel rigidity]\label{thm:one-channel}
Let $s=B+1$. The infimum of $\nustd$ over finite affine lifts through
$\mathscr K$ equals one if the dictionary contains $Q_{B+2}$, and equals
two otherwise. In the latter case the corresponding infimum for
$(\B_2^{B+1})^b$ equals $2b$. Moreover, at some simultaneous contact
there are pure-row certificates, one per ball, whose positive
aggregates all have total Jordan rank at least $2b$.
\end{theorem}
\begin{proof}
Suppose first that one lift has $\nustd<2$. The integral boundary orders
in Lemma~\ref{lem:boundary-order} imply that every boundary tuple has
nullity at most one. Thus every member of every normalized genuine
certificate fiber $\mathcal D(v)$ has rank at most one. Such a member is
nonzero, since its slack is not identically zero. The positive midpoint
of two elements on different product-cone rays has rank at least two.
Convexity therefore puts $\mathcal D(v)$ on one ray. The whole-slice
identity fixes its scalar, making $\mathcal D(v)=\{Y(v)\}$.

A Slater point provides the uniform normalization
\[
 0<1-\norm{\pi X^\circ}\leq\ip{X^\circ}{Y(v)}
 \leq1+\norm{\pi X^\circ}.
\]
Strict positivity of $X^\circ$ bounds $Y(v)$ above and away from zero.
Since $Y$ is bounded with closed graph, it is continuous. As the
sphere is connected, the factor in which $Y(v)$ is nonzero does not
depend on $v$. Projectivization gives a continuous injection
$\Sph^B\to\mathcal P(V_j)$ into that factor's primitive-ray manifold:
if the rays agree, evaluating the genuine identities at a lift of the
ball's center fixes their scalar, and the identities then fix $v$.

The target dimension is $a_j(r_j-1)\leq B$. Invariance of dimension
excludes a smaller dimension, and invariance of domain makes the map a
homeomorphism in equal dimension. The primitive-ray spaces are
$\mathbb{RP}^{r-1}$, $\mathbb{CP}^{r-1}$,
$\mathbb{HP}^{r-1}$, spheres for spin factors, and $\mathbb{OP}^2$ for
the Albert algebra. The classification and cohomology used in
Section~\ref{sec:orbits} show that the only sphere cases are the spin
factors, including the order-two Hermitian cones, which are Lorentz
cones. Thus a value below two requires $Q_{B+2}$ in the dictionary. Its
direct section
attains one; otherwise the two-group construction
\eqref{eq:fixed-tail} attains two.

For the product assertion apply sequential compression as in
Section~\ref{sec:products}. At the next source, compress the entire
original certificate fiber to the faces complementary to the previously
chosen support join. The resulting families are compact, convex, and
have closed graphs: the original fibers are uniformly bounded by the
same Slater normalization, and compression is linear. Their genuine
row identities on the remaining contact cylinder ensure that
zero is not in any compressed fiber. If every compressed certificate had
total rank at most one, the preceding singleton and primitive-ray
argument would again give a spin factor of capacity $B$ in one of the
remaining faces. A proper face of a simple factor has strictly smaller
primitive capacity than that factor; hence no such face can arise when
the original dictionary does not contain $Q_{B+2}$. Some compressed
certificate must therefore have rank at least two. The support-join
compression identity of Section~\ref{sec:products} adds at least two
new rank units at each source. The final join, and every positive
aggregate of the chosen certificates, has rank at least $2b$.
All tuples at that simultaneous contact are complementary to this join,
so the boundary-order test gives $\nustd\geq2b$. Separate two-group
constructions attain equality.
\end{proof}

For example, an Albert-cone dictionary has $B=16$, and the minimum of
$\nustd$ for the $17$-dimensional ball is two. The theorem concerns the
optimum over lifts; it does not say that every lift using a matching
spin factor has parameter one.

\subsection{Arbitrarily many factors under a wide cap}
For Hermitian order caps and Lorentz dimension caps, we settle the case
$s-1=qB$ with $q\geq2$ and $B\geq q-1$, again without selection
hypotheses or a bound on the number of factors.

\begin{lemma}[A face-codimension budget]\label{lem:face-codim}
Suppose a bounded relative-Slater affine slice $P=L\cap K$ projects onto
$\B_2^s$. For any $X\in P$ over the boundary of the ball, let $E_X$ be the linear
span of the minimal face of $K$ containing $X$. Then
$\codim E_X\geq s$.
\end{lemma}
\begin{proof}
Write $L=X^\circ+T$. Since $P$ is bounded and nonempty,
$T\cap K=\{0\}$; also $0\notin L$, since otherwise $P$ would be a
bounded cone with nonempty relative interior projecting onto a ball.
On $\widehat L=\spanop\{X^\circ\}+T$ define height
$\ell(tX^\circ+H)=t$. A cone point of zero height is zero.
There are no cone points of negative height: otherwise adding a suitable
positive multiple of $X^\circ$ gives a nonzero cone point of height zero;
pointedness excludes cancellation to zero. Hence
$\widehat L\cap K$ is precisely the cone over $P$.
The projection homogenizes to a linear surjection
$\widehat\pi:\widehat L\to\R^{s+1}$ whose cone image is $Q_{s+1}$.
If $k=\dim\ker\widehat\pi$, then $\dim\widehat L=s+1+k$.
Every $H\in\widehat L\cap E_X$ permits two-sided small feasible
perturbations of $X$, because $X\in\ri(K\cap E_X)$. Their images lie
in the line spanned by the target extreme ray. Thus
$\dim(\widehat L\cap E_X)\leq k+1$.
Grassmann's dimension inequality gives
$k+1\geq s+1+k-\codim E_X$, as required.
\end{proof}

\begin{theorem}[Wide-cap rigidity]\label{thm:wide-cap}
Fix one of two dictionaries: Hermitian cones over one field
$\mathbb F\in\{\R,\C,\HH\}$ of order at most $R$, with $B=a(R-1)$;
or Lorentz cones of dimension at most $d$, with $B=d-2$. Consider
lifts through any finite number of rays and of factors from this
dictionary. If $q\geq2$ and
$B\geq q-1$, the minimum of $\nustd$ for affine lifts of
$\B_2^{qB+1}$ is $q+1$.
\end{theorem}
\begin{proof}
Unbounded lifts were settled by Proposition~\ref{prop:barrier-frontiers}.
For a bounded lift suppose $\nustd<q+1$. All boundary tuples then have
integer total nullity at most $q$. In a Hermitian block of order $r\leq R$
and nullity $m$, the minimal-face codimension is
\[
 \left[r+\frac{ar(r-1)}2\right]
 -\left[r-m+\frac{a(r-m)(r-m-1)}2\right]
 =m+\frac{am(2r-m-1)}2\leq(B+1)m.
\]
A Lorentz block has codimension $d_i-1\leq B+1$ at a nonzero boundary
point, nullity one there, and codimension $d_i\leq2(B+1)$ at its
vertex, nullity two there. Rays obey the same bound. Hence
$\codim E_X\leq(B+1)\operatorname{nul}X$.
Lemma~\ref{lem:face-codim} rules out nullity at most $q-1$, since
$(q-1)(B+1)<qB+1$ under the stated hypothesis. Every tuple in every
boundary fiber therefore has nullity exactly $q$.

Each compact boundary fiber is a singleton. Indeed, from a
relative-interior point of a nontrivial fiber take a maximal nonzero
affine chord. Its endpoint must leave the relative interior of the
minimal cone face, and thus acquire an additional null eigenvalue.
The resulting nullity exceeds $q$, a contradiction. Compactness makes
the unique boundary lift $X(v)$ continuous.

With $p_i$ and $q_i$ the Jordan ranks of complementary primal and
certificate blocks, as in \eqref{eq:rank-lower-chain}, the curvature
chain holds with equality on a dense open set of supports:
\[
 qB\leq\sum_i a_ip_iq_i\leq
 \sum_i a_i(r_i-q_i)q_i\leq B\sum_iq_i\leq qB.
\]
Hence exactly $q$ cap-attaining blocks have corank one, and all other
blocks are interior. Call the set of corank-one blocks at a generic
point its active pattern. At a limit of points with one pattern, those
zero eigenvalues persist, and constant total nullity excludes any
additional zero eigenvalue or new active block. The closures of the
regions with distinct patterns are therefore disjoint, and their finite
union covers the connected sphere, so one pattern occurs everywhere.

The corresponding kernel lines define a continuous map from
$\Sph^{qB}$ to $(\mathbb FP^{R-1})^q$, or to $(\Sph^B)^q$ in the
Lorentz case. It is injective. Indeed, equal lines put the two tuples in
the relative interior of the same product face. Any certificate exposing
one projected point annihilates that face, so the genuine slack identity
makes the other projected point identical. Source and target are compact
connected manifolds of the same dimension, so invariance of domain would
make the map a homeomorphism. The target has nonzero intermediate
mod-two cohomology, whereas the sphere does not. This contradiction
proves the lower bound. The grouped construction with $q+1$ groups
attains it.
\end{proof}

The proof also gives a statement that does not depend on the cap width:
a compact lift with $s-1=qB$, $q\geq2$, cannot have total nullity exactly
$q$ at every point of every boundary fiber. Such constant nullity would
force singleton fibers and hence the same impossible kernel-line
injection.

\subsection{Properties of a hypothetical better bounded lift}
For a fixed Hermitian order cap or Lorentz dimension cap, the single-ball
cases not resolved by the preceding results are those with
\begin{equation}\label{eq:unsettled-range}
 s=qB+1,\qquad q\geq B+2,
 \qquad q\leq\inf_{\text{affine lifts}}\nustd\leq q+1.
\end{equation}
By Proposition~\ref{prop:barrier-frontiers}, any lift with $\nustd<q+1$
must be bounded. The next proposition derives necessary properties of
such a lift; none of the preceding theorems depends on resolving this
interval.

\begin{proposition}[Certificate collapse and fixed generic labels]
\label{prop:range-collapse}
Fix a field $\mathbb F\in\{\R,\C,\HH\}$, an order cap $R\geq2$,
and $B=a(R-1)$, where $a=\dim_{\R}\mathbb F$. Consider a bounded
affine lift through finitely many Hermitian cones of orders at most $R$
and optional rays, in the divisible case $s-1=qB$, $q\geq2$.
Suppose $\nustd<q+1$. For $v\in\Sph^{qB}$ let
$\mathcal P(v)$ and $\mathcal D(v)$ be its full primal and normalized
certificate fibers. Define
\[
 E_i(v)=\sum_{Y\in\mathcal D(v)}\operatorname{Ran}_{\mathbb F}Y_i,
 \qquad Z=\left\{v:\sum_i\dim_{\mathbb F}E_i(v)\leq q-1\right\},
\]
counting a nonzero ray certificate as dimension one. Then $Z$ is nonempty
and semialgebraic, and $\dim Z\leq qB-B$. On a dense open semialgebraic
set $U$, both fibers are singletons, strict complementarity holds, and
exactly $q$ full-order blocks are of corank one.
If $B\geq2$, one and the same set of $q$ block labels is active on all
of $U$. Thus changing generic label patterns cannot produce such a lift;
what we cannot exclude is that this fixed pattern degenerates over $Z$.
\end{proposition}
\begin{proof}
The boundary-order bound gives nullity at most $q$ at every primal
boundary tuple, so every certificate has rank at most $q$. A finite
positive average spans all the spaces $E_i(v)$, hence
$\sum_i\dim E_i(v)\leq q$. Semialgebraicity follows by expressing the
spanning condition with finitely many certificates (the ambient dimension
bounds their number) and eliminating quantifiers.

Choose a dense open generic set where the full-fiber minimum-rank theorem
holds. All certificates there have rank exactly $q$, and all primal
tuples have nullity exactly $q$. The compact primal chord argument in
Theorem~\ref{thm:wide-cap} gives a singleton primal fiber. The dual fiber
lies in a product face of rank at most $q$; if it contained a nontrivial
maximal chord, an endpoint would leave that face's relative interior and
have rank below $q$. Thus it too is a singleton. On a further dense open
smooth locus, equality in the curvature chain gives exactly $q$
full-order corank-one blocks and strict complementarity.

On each smooth stratum of $Z$ of dimension $t$, choose smooth primal and
certificate selections after further stratification. Differentiating the
slack on two copies of this stratum gives its positive definite tangent
metric, of rank $t$. The certificate ranks sum to at most $q-1$, so the
mixed Peirce estimate gives $t\leq B(q-1)$. This proves the dimension
bound, including for the semialgebraic closure $\overline Z$.

For any support $v$, compactness gives a limit of generic primal tuples.
Its nullity is at least $q$, because nullity cannot drop in a limit, and at
most $q$ by the parameter hypothesis. If $v\notin Z$, an averaged
certificate of rank $q$ forces every primal in that fiber to have
nullity $q$. The fiber is consequently a singleton. This singleton is
the limit of every approaching sequence of generic tuples. Passing to
a subsequence of fixed generic labels shows that it has exactly that
pattern: an additional zero eigenvalue would exceed nullity $q$.
Its inactive blocks are positive definite. Compactness then provides a
neighborhood of $v$ in which no generic tuple has a new active label;
since each generic tuple has exactly $q$ labels, its pattern is fixed.

If $Z$ were empty, this argument would give a continuous singleton lift
and one fixed pattern on the whole sphere. The kernel-line injection
and projective-product contradiction in Theorem~\ref{thm:wide-cap}
would apply. Hence $Z\ne\varnothing$.
If $B\geq2$, the complement of the closed semialgebraic set
$\overline Z$, of codimension at least two, is connected. For example,
semialgebraic triangulation and a small general-position perturbation
of a path avoid a subcomplex of codimension at least two. The local
pattern just proved is constant on this complement. Every nonempty
open generic pattern region meets the complement because
$\overline Z$ has empty interior. There can therefore be only one
such pattern.
\end{proof}

For Lorentz cones of dimension at most $d\geq3$, set $B=d-2$ and
retain the same boundedness and parameter hypotheses, with optional rays.
The Lorentz version replaces $E_i$ by the support face generated by the
whole certificate fiber, and its dimension count by Jordan rank. Its
rank-one orbit has dimension at most $B$. The same proof gives
$\dim Z\leq qB-B$ and, for $B\geq2$, the same fixed-label conclusion.
When $B=1$, the exceptional set may have codimension one and the
connected-complement argument does not apply.

We record one precise topological fact about the exceptional fibers.
The contact incidence space
\[
 \Gamma=\{(v,X,Y):v\in\Sph^{qB},\ X\in\mathcal P(v),\
 Y\in\mathcal D(v)\}
\]
is compact: primal compactness is assumed and the Slater normalization
uniformly bounds all dual fibers. The fiber of $\Gamma\to\Sph^{qB}$
is exactly $\mathcal P(v)\times\mathcal D(v)$, since every primal is
complementary to every genuine certificate. This is a nonempty compact
convex set. The Vietoris--Begle theorem for compact spaces with acyclic
fibers \citep{Begle1950} implies
\begin{equation}\label{eq:incidence-cohomology}
 \check H^k(\Gamma;\mathbb Z/2)
 \cong\check H^k(\Sph^{qB};\mathbb Z/2).
\end{equation}
Thus the incidence space has no cohomology beyond that of the sphere.
The difficulty is that the generic kernel-line map need not extend over
the rank-deficient fibers. Equation~\eqref{eq:incidence-cohomology}
gives neither such an extension nor its degree, and we do not claim
that one exists.

\subsection{Counterexamples to stronger fiberwise conclusions}

\begin{example}[Boundary selections can escape over a compact target]
\label{ex:escaping-disk}
Put $\delta=1-x_1$ and impose three positive semidefinite blocks:
\[
 \begin{pmatrix}1+x_1&x_2\\x_2&\delta\end{pmatrix}\succeq0,
 \qquad
 \begin{pmatrix}\delta&x_2\\x_2&u\end{pmatrix}\succeq0,
 \qquad
 \begin{pmatrix}\delta&u\\u&z-1\end{pmatrix}\succeq0.
\]
The first block describes exactly $\B_2^2$. For every point of this
ball with $\delta>0$, the remaining blocks are feasible precisely when
$u\geq x_2^2/\delta$ and $z\geq1+u^2/\delta$. At $(x_1,x_2)=(1,0)$,
they instead require $u=0$ and $z\geq1$. Thus the projection is the
entire closed disk. The choice $x=0$, $u=1$, $z=3$ makes all three
blocks positive definite, so the lift has Slater feasibility.
On the boundary $x_1=1-\delta$, $x_2^2=2\delta-\delta^2$, every
feasible completion satisfies
\[
 u\geq2-\delta,\qquad
 z\geq1+\frac{(2-\delta)^2}{\delta}\longrightarrow\infty
 \quad(\delta\downarrow0).
\]
Consequently no boundary selection is locally bounded at $(1,0)$.
Compactness of the projected set and Slater feasibility therefore do
not imply bounded boundary selections; any compactness assumption on the lifted
boundary needs separate justification.
\end{example}

\begin{lemma}[A compact base removes one gradient unit]
\label{lem:compact-base-parameter}
Let $F$ be a $k$-logarithmically homogeneous self-concordant barrier on a
pointed cone, with positive definite Hessian on its span. On a nonempty
base $\ell(X)=1$, where $\ell$ is nonnegative on the cone, its restricted
gradient parameter is at most $k-1$.
\end{lemma}
\begin{proof}
Write $g=\nabla F$, $H=\nabla^2F$. Homogeneity gives
$HX=-g$ and $X^THX=k$. The squared dual norm of the gradient on
$\ker\ell$ is its orthogonal projection in the inverse-Hessian norm:
\[
 \min_\alpha(g-\alpha\ell)^TH^{-1}(g-\alpha\ell)
 =k-\frac1{\ell^TH^{-1}\ell}.
\]
The unit Dikin ellipsoid lies in the cone; since $\ell(X)=1$ and
$\ell$ is nonnegative there, $\ell^TH^{-1}\ell\leq1$. This proves the
claim. Compactness of the base is a convenient sufficient condition
for a strictly positive height; the argument only needs the conditions
stated above.
\end{proof}

\begin{example}[Small minimum nullity, larger global parameter]
\label{ex:rotated-nullity}
The Lorentz case of the perspective construction in
Section~\ref{sec:rank} shows that the standard barrier parameter can
exceed both the minimum fiber nullity and the maximum certificate rank;
all three are computed exactly. Partition
$\R^{s-1}$ into $q$ nonempty coordinate groups and write $x=(a,b)$.
The lift is
\begin{equation}\label{eq:rotated-counterexample}
 W_G=\left(\frac{1-b+z_G}2,\frac{1-b-z_G}2,a_G\right)\in Q_{|G|+2},
 \qquad\sum_Gz_G=1+b.
\end{equation}
Equivalently, $1-b\geq0$, $z_G\geq0$, and
$\norm{a_G}^2\leq(1-b)z_G$. Summing them shows that the projection lies
in the ball. Conversely, every ball point lifts by distributing the nonnegative residual $1-b^2-\norm a^2$
among the $z_G$. All variables are bounded and $a=b=0$, $z_G=1/q$
is strictly feasible.

At every boundary point with $b<1$, every inequality is an equality,
$z_G=\norm{a_G}^2/(1-b)$, and each block is a nonzero extreme ray.
At the north pole, the fiber is $z_G\geq0$, $\sum_Gz_G=2$.
Its relative interior again has one null eigenvalue per block, whereas
a vertex has one such block and $q-1$ zero blocks. Consequently
\[
 \min_{X\in\mathcal P(v)}\operatorname{nul}X=q
 \quad\hbox{for every }v,\qquad
 \max_{X\in\mathcal P(\text{north})}\operatorname{nul}X=2q-1.
\]
For completeness, at support $v=(c,\eta)$ the entire certificate fiber
has blocks
\[
 V_G=(A_G+\gamma,A_G-\gamma,-c_G),\quad
 \gamma=(1-\eta)/2,\quad\sum_GA_G=(1+\eta)/2,\quad
 A_G\geq0,\quad 4A_G\gamma\geq\norm{c_G}^2.
\]
Coefficient matching on the entire affine slice proves these conditions
necessary and sufficient. For $\eta<1$, all the inequalities are equalities
and the certificate is unique with every block nonzero of rank one.
At $\eta=1$ it is the simplex $A_G\geq0$, $\sum_GA_G=1$.
Thus the maximum certificate rank is $q$ at every support, although its
minimum at the north pole is one.

The restricted standard barrier is
$F=-\sum_G\log((1-b)z_G-\norm{a_G}^2)$, with exact parameter $2q-1$.
For the upper bound homogenize $1$ to $\tau$ in
\eqref{eq:rotated-counterexample}; the resulting cone has a compact
$\tau=1$ base and its restricted product barrier is $2q$-homogeneous.
Apply Lemma~\ref{lem:compact-base-parameter}. The north-fiber vertex
and Lemma~\ref{lem:boundary-order} give the matching lower bound.
For balls with different $q_a$, separate copies of this lift give
minimum fiber nullity $\sum_aq_a$, maximum rank $\sum_aq_a$ of aggregate
certificates with strictly positive source weights, and standard
parameter $\sum_a(2q_a-1)$. Constant terms cannot move between source
certificates: a source affine slack
$\mu_a-\lambda_av_a^Tx_a$ is nonnegative only if
$\mu_a\geq\lambda_a$, and equality of the total constants forces
equality sourcewise.
\end{example}

This example refutes an extra-unit lower bound stated with the
quantifier ``every primal completion at some support'': every support
has a completion of nullity only $q$. The barrier lower bound of
Lemma~\ref{lem:boundary-order} needs only one completion of high
nullity, such as the north-fiber vertex.

Two simpler examples show that convexity of the fibers alone does not
yield a global kernel-line map into a product of projective spaces. In a two-dimensional $\mathbb F$-subspace
$U$, the set $\{Y\succeq0:\operatorname{Ran}Y\subset U,\tr Y=1\}$
contains all primitive rays in $U$, normalized to trace one, and their
rank-two mixtures. At its
scalar center there is no distinguished eigenline. Adding $q-2$ fixed
rank-one blocks keeps the total rank at most $q$.
For a primal example, encode the simplex
$\sum_{j=0}^q p_j=1$, $p_j\geq0$, using
\[
 X_0=\operatorname{diag}(p_0,p_1),\qquad
 X_j=\operatorname{diag}(1,p_{j+1}),\quad1\leq j\leq q-1,
\]
and project to $\sigma=\sum_{j=2}^qp_j\in[0,1]$.
The fiber at $\sigma=0$ has $p_0=t$, $p_1=1-t$: its first kernel
changes between two lines at the endpoints and disappears in between.
Its nullity is $q-1$ inside the fiber and $q$ at either endpoint.
The standard barrier is $-\sum_{j=0}^q\log p_j$, with exact parameter
$q$ by Lemmas~\ref{lem:compact-base-parameter} and
\ref{lem:boundary-order}. This is a lift of an interval, not of a ball
in the critical dimension. It shows that a proof for the unresolved
range \eqref{eq:unsettled-range} must use more than boundedness,
convexity, and a nullity bound.
